\documentclass{article}
\usepackage[utf8]{inputenc}
\usepackage[T1]{fontenc}
\usepackage{graphicx}
\usepackage{algorithm}
\usepackage{algpseudocode}

\usepackage{listings}
\usepackage{xcolor}
\usepackage[margin=2cm]{geometry}
\usepackage{float}

\usepackage{amssymb} %Ajoute des symboles mathématiques comme R, Z, forall, exists, etc....
\usepackage{amsmath}
\usepackage{amsthm}

\usepackage{enumitem}


\lstdefinestyle{monstyle}{
    backgroundcolor=\color{white},
    commentstyle=\color{gray},
    keywordstyle=\color{blue},
    numberstyle=\tiny\color{gray},
    basicstyle=\ttfamily\small,
    breaklines=true,
    numbers=left,
    numbersep=5pt,
    frame=single,
}
\lstset{style=monstyle}


\usepackage{soul} % IMPORTANT : POUR SURLIGNER AVEC \hl en fin de ligne 
\usepackage{xcolor}


\title{\sethlcolor{purple!30}\hl{\textbf{Fiche de révisions: SSH -- Environnement de travail}}}
\author{Derrez Lorenzo}

\begin{document}
\maketitle

\section*{\hl{\textbf{Introduction}}}

\begin{itemize}
    \item Permet d'ouvrir un terminal sur une machine à distance
    \item Chiffrée et authentifiée
    \item Travailler à distance dans un shell 
    \item exécuter une commande sur une autre machine 
    \item Copier des fichiers avec scp ou sftp 
    \item monter un système de fichiers distant avec sshfs
    \item \textbf{Remarque:} une fois connecté, les commandes agissent sur la machine distante, pas sur la machine locale.
\end{itemize}

\section*{\hl{\textbf{Fonctionnement}}}

\subsection*{\textbf{Connexion à une machine distante:}}

\textbf{ssh [-X] login@hostname}

\begin{itemize}
    \item login : votre login sur la machine distante; 
    \item hostname : nom de la machine distante; 
    \item -X : redirection graphique si elle est nécessaire 
    \item vous arrivez dans votre HOME distant
\end{itemize}

\section*{\hl{\textbf{Copies à distance}}}

\subsection*{\textbf{Copie \underline{vers} une machine distante:}}
\textbf{scp [-r] src1 [src2 ... srcN] login@hostname:dst\_path}

\begin{itemize}
    \item dst\_path : chemin absolu ou relatif au répertoire personnel distant; 
    \item -r : copie récursive d'un répertoire; 
    \item les règles sont proches de celles de cp.
\end{itemize}

\subsection*{\textbf{Copier \underline{depuis} une machine distante:}}

\textbf{scp login@hostname:src\_path dst\_ path}

\section*{\hl{\textbf{Les clés SSH}}}

\begin{itemize}
    \item il est lourd de devoir retaper son mot de passe ou sa passphrase à chaque connexion/execution 
    \item Il est possible d'utiliser un mécanisme basé sur un système de clé privé/clé publique; 
\end{itemize}

\textbf{Principe:}
\begin{itemize}
    \item Clé publique à copier sur la machine distante: peut être vue comme un verrou.
    \item Clé privée qui \textbf{reste sur la machine locale}: peut être vu comme l'unique clé qui ouvre le verrou.
    \item \textbf{REGLE :} ne jamais copier, envoyer, ou partager la clé privée.
\end{itemize}

\textbf{Créer un couple clés privée/publique:}
\begin{enumerate}
    \item Créer une clé avec ssh-keygen : 
    \subitem ssh-keygen -t ed25519
    \item Accepter le nom proposé par défaut pour le fichier de clé.
    \item Mettre une passphrase : une phrase qui protège la clé.
    \item Vérifier que deux fichiers ont été crées.
\end{enumerate}

\end{document}